\scanpage{222}\mspage{167}

% MS 1–2
C'est nécessaire, car \deleted{\unsure{par le}} si $\unsure{\xi|\zeta}=\mathrm{hot}_{\unsure{W|w}}(\unsure{T|I})$, \deleted{donc} donc $\xi_{\unsure{S+|S}}=\mathrm{hot}_{S,W}(T_S)$ (où $T_S=T\times S\,/\,S$)

% MS 3–4
\deleted{donc $T_S$ parfait sur $S$, et par suite ($T_S$ étant lisse sur $S$)}

% MS 5
\[
 \deleted{\xi_{X'}=p'^*(T_S),\quad \xi_X=p^*(T_S)}
\]

% MS 6
et on a donc

% MS 6–7
\[
\begin{aligned}
 &\deleted{\unsure{p'_!}(\xi_{X'})=\unsure{\mathrm{hot}_{S'}}}\\
 &\deleted{p'_!p'^*(T_S)\simeq\unsure{X'}\unsure{;}}
\end{aligned}
\]

% MS 8
\[
 p'_!(\xi_{X'})=\mathrm{hot}_{S,W}(X'\times T),\qquad
 p_!(\xi_X)=\mathrm{hot}_{S,W}(X\times T),
\]

% MS 9–11
et il faut prouver que si $u:X'\longrightarrow X$ est dans $W_S^d$, alors $\underbrace{X'\times T}_{X'\times_S T_S}\longrightarrow\underbrace{X\times T}_{X\times_S T_S}$ aussi, ce qui \uline{signifie} simplement que \deleted{$({}_{s\backslash}\unsure{X'}\times T)\longrightarrow({}_{s\backslash}X)\times T$}

% MS 12, label at left
\[
 \deleted{\unsure{(}}{}_{s\backslash}(X'\deleted{\unsure{\xi}}\times T)\longrightarrow{}_{s\backslash}(X\times T)\tag*{\unsure{les|(*)}}
\]

% MS 13–17
est dans $W$ pour \unsure{tt} $s\in S$. Mais le changement de base $ {}_{s\backslash}S\longrightarrow S$ commute aux produits fibrés, donc transforme $X\times T=X\times_S T_S$ en $({}_{s\backslash}X)\times_{{}_{s\backslash}S}T_{({}_{s\backslash}S)}={}_{s\backslash}X\times T$, et de $=$ pour $X'$. Donc la flèche s'identifie à

% MS 18
\[
 ({}_{s\backslash}X')\times T\longrightarrow({}_{s\backslash}X)\times T
\]

% MS 19–23
déduite de $ {}_{s\backslash}X'\longrightarrow{}_{s\backslash}X$ par \deleted{\unsure{un bifoncteur}} le foncteur $?\times T$. Comme ce foncteur transforme $W$\unsure{-équiv|-éq} en $W$, on gagne. \deleted{Finalement, le cas b) du \unsure{théorème} du th.~4 se présente ainsi :}

% archive sheet number: 221
